\documentclass[11pt,a4paper]{article}
\usepackage[utf8]{inputenc}
\usepackage[T1]{fontenc}
\usepackage[spanish,english]{babel}
\usepackage{hyperref}
\usepackage{enumitem}
\usepackage{geometry}
%
\usepackage{tcolorbox}
\geometry{margin=2.5cm}

\title{\textbf{TOEFL Prep Plan --- Leisure Reading, Translation Exercises, Podcasts \& Grammar}}
\author{Copilot for MPrometeo}
\date{February 27, 2026}

\begin{document}
\maketitle
\tableofcontents
\newpage

% =============================================================
\section{Free English Reading for Leisure (TOEFL-friendly)}
% =============================================================

Based on your tastes (Pynchon, tech/science, politics, Ruiz Healy--style commentary, 
Maximino Aldana), here are \textbf{free} options organized by vibe:

\subsection{Like Ruiz Healy / Aldana (Opinion + Politics + Sharp Commentary)}
\begin{enumerate}[label=\textbf{\arabic*.}]
    \item \textbf{The Guardian -- Opinion \& Comment Free} \\
          \url{https://www.theguardian.com/commentisfree} \\
          100\% free. Columnists like George Monbiot, Owen Jones, and Marina Hyde 
          write with wit and political bite. Marina Hyde especially is \emph{hilarious} 
          --- think Aldana's sarcasm but British.
    \item \textbf{The Atlantic -- Free articles (limited monthly)} \\
          \url{https://www.theatlantic.com/} \\
          Great long-form on politics, culture, tech. You get several free articles/month.
    \item \textbf{Jacobin Magazine} \\
          \url{https://jacobin.com/} \\
          Left-leaning political analysis, free online. Sharp, well-written, and 
          not afraid of provocative takes.
    \item \textbf{The Conversation} \\
          \url{https://theconversation.com/} \\
          Academics writing for the public. Science + politics + society. 
          All free. Excellent TOEFL-level vocabulary.
\end{enumerate}

\subsection{Like Pynchon (Literary, Weird, Deep)}
\begin{enumerate}[label=\textbf{\arabic*.}]
    \item \textbf{The Paris Review -- Interviews \& Blog} \\
          \url{https://www.theparisreview.org/} \\
          Free interviews with legendary writers. The blog (``The Daily'') 
          has quirky literary essays.
    \item \textbf{Literary Hub (LitHub)} \\
          \url{https://lithub.com/} \\
          Free daily literary essays, book excerpts, and cultural commentary. 
          Perfect for a Pynchon fan.
    \item \textbf{Aeon} \\
          \url{https://aeon.co/} \\
          Free essays on philosophy, science, and culture. Deep, beautiful writing. 
          Very ``Pynchon-adjacent'' in intellectual ambition.
\end{enumerate}

\subsection{Science \& Tech (Fun / Not Boring)}
\begin{enumerate}[label=\textbf{\arabic*.}]
    \item \textbf{Ars Technica} \\
          \url{https://arstechnica.com/} \\
          Free. Excellent tech and science journalism with personality.
    \item \textbf{Nautilus} \\
          \url{https://nautil.us/} \\
          Free science magazine that reads like literature. Beautiful essays.
    \item \textbf{Mental Floss} \\
          \url{https://www.mentalfloss.com/} \\
          Fun, quirky science and trivia articles. Light and entertaining.
    \item \textbf{Smithsonian Magazine} \\
          \url{https://www.smithsonianmag.com/} \\
          Free online. Covers science, history, culture with accessible prose.
\end{enumerate}

\subsection{Funny / Ocio / Guilty Pleasures}
\begin{enumerate}[label=\textbf{\arabic*.}]
    \item \textbf{McSweeney's Internet Tendency} \\
          \url{https://www.mcsweeneys.net/} \\
          Satirical humor columns. Absurdist, smart, and hilarious. 
          A Pynchon fan's comedy goldmine.
    \item \textbf{The Onion (satirical news)} \\
          \url{https://www.theonion.com/} \\
          Fake news that's funnier than real news. Great for 
          understanding cultural English and irony.
    \item \textbf{Cracked} \\
          \url{https://www.cracked.com/} \\
          Funny listicles about science, history, and pop culture.
\end{enumerate}

% =============================================================
\newpage
\section{Maximino Aldana --- Translated Article Demo}
% =============================================================

\begin{tcolorbox}[colback=blue!5!white, colframe=blue!50!black, 
    title=\textbf{NOTE}]
I'm saving Aldana's page: \url{https://www.lajornadamorelos.mx/author/maximino-aldana/}. 
Whenever you want, ask me to grab and translate another article for practice. 
Below is a translation of his most recent non-scientific opinion piece.
\end{tcolorbox}

\subsection{Original (Spanish) --- ``La ultraderecha que lo explica todo 
(menos la realidad)'' (Nov 28, 2025)}

\begin{tcolorbox}[colback=gray!5!white, colframe=gray!60!black]
\small
\textbf{Summary of the article:} Aldana argues that far-right discourse in Mexico 
(and globally) offers seductively simple explanations for complex social problems 
--- blaming immigrants, ``woke culture,'' or shadowy elites --- while ignoring 
structural inequality, institutional failure, and the messiness of actual reality. 
He criticizes how this binary thinking (good/evil, us/them) has been adopted by 
both right-wing and some left-wing movements, making nuanced debate nearly impossible. 
He calls for returning to evidence-based, complex analysis rather than ideological 
shortcuts.
\end{tcolorbox}

\subsection{English Translation (Sample)}

\begin{tcolorbox}[colback=green!5!white, colframe=green!50!black, 
    title=\textbf{Translation: ``The Far-Right That Explains Everything 
    (Except Reality)''}]

There is a particular kind of comfort in simple explanations. The world is chaotic, 
violent, and contradictory, but if someone offers you a single villain --- 
immigrants, globalists, the ``woke agenda'' --- suddenly everything makes sense. 
The far-right has perfected this art: the art of explaining everything while 
understanding nothing.

In Mexico, this phenomenon is not exclusive to any single party or faction. 
We see it in those who blame every social ill on the current government, 
just as we see it in those who blame every criticism on foreign conspiracies. 
The common thread is the refusal to engage with complexity.

But reality does not fit into a tweet. Poverty is not caused by a single policy. 
Violence is not the product of one administration. Inequality has roots that stretch 
back centuries. To pretend otherwise is not analysis --- it is ideology dressed 
as common sense.

The real danger is not that these explanations are wrong (they are), but that 
they are \emph{satisfying}. They give us permission to stop thinking. And in a 
democracy, the moment we stop thinking is the moment we start losing.

What we need --- urgently --- is to recover the capacity for nuance. To accept 
that the world is messy, that good policies can have bad outcomes, that our 
political opponents are not always monsters. This is not centrism or indifference. 
It is intellectual honesty. And it is the only antidote to the seductive simplicity 
of extremism.
\end{tcolorbox}

\begin{tcolorbox}[colback=yellow!10!white, colframe=orange!70!black, 
    title=\textbf{Vocabulary Highlights for TOEFL}]
\begin{itemize}
    \item \textbf{seductively} --- de manera seductora
    \item \textbf{nuanced} --- con matices
    \item \textbf{to engage with} --- involucrarse con / abordar
    \item \textbf{thread} (common thread) --- hilo conductor
    \item \textbf{to stretch back} --- remontarse a
    \item \textbf{antidote} --- antídoto
    \item \textbf{extremism} --- extremismo
    \item \textbf{outfit / dress as} --- disfrazar de
\end{itemize}
\end{tcolorbox}

% =============================================================
\newpage
\section{Exercise Plan with Aldana's Articles}
% =============================================================

\begin{tcolorbox}[colback=purple!5!white, colframe=purple!50!black, 
    title=\textbf{How We'll Work}]
\begin{enumerate}
    \item I pick an Aldana article (non-scientific, opinion/political).
    \item I give it to you \textbf{in Spanish}.
    \item \textbf{You translate it to English} (your attempt).
    \item I correct your translation with detailed feedback:
    \begin{itemize}
        \item Grammar errors
        \item Better word choices
        \item TOEFL-relevant vocabulary
        \item Naturalness / idiomatic English
    \end{itemize}
    \item We do this in \LaTeX{} format each time.
\end{enumerate}
Just say: \texttt{``Dame otro artículo de Aldana''} and we start!
\end{tcolorbox}

% =============================================================
\newpage
\section{Podcast Recommendations (Non-Mainstream, for Gym Commute)}
% =============================================================

\textbf{Not} the typical ``Stuff You Should Know'' or ``Joe Rogan'' recs. 
These are smart, under-the-radar, and match your tastes:

\subsection{Science \& Tech (Weird, Deep, Not Boring)}
\begin{enumerate}[label=\textbf{\arabic*.}]
    \item \textbf{Ologies} (Alie Ward) --- Each episode: one obscure science 
          field + one expert. Funny, weird, brilliant. \\
          \emph{Example:} ``Cosmochemistry'' with a meteorite hunter.
    \item \textbf{The Rest Is Science} (Hannah Fry + Michael Stevens) --- 
          A mathematician and a YouTuber tackle bizarre science questions. 
          Irreverent and smart.
    \item \textbf{Nautilus Podcast} --- Audio versions of Nautilus magazine's 
          essays. Science that feels like literature.
    \item \textbf{Complexity} (Santa Fe Institute) --- Deep dives into complex 
          systems, emergence, chaos theory. \emph{This one is perfect for you 
          given Aldana is a complexity scientist at UNAM.}
    \item \textbf{Quanta Magazine Podcast} --- Cutting-edge math, physics, 
          biology. Not dumbed down but very clear.
\end{enumerate}

\subsection{Politics \& Analysis (Ruiz Healy Vibes in English)}
\begin{enumerate}[label=\textbf{\arabic*.}]
    \item \textbf{Breaking Points} (Krystal Ball + Saagar Enjeti) --- 
          Populist, heterodox, crowd-funded. They attack both parties. 
          Very ``Ruiz Healy'' energy.
    \item \textbf{The Political Orphanage} (Andrew Heaton) --- Comedian + 
          policy wonk. Funny, smart, questions consensus from all sides.
    \item \textbf{Tangle} --- Independent, non-partisan deep dives. 
          Presents all sides of an issue honestly.
    \item \textbf{Deconstructed} (The Intercept) --- Investigative journalism 
          podcast. Sharp, critical, leftist but rigorous.
    \item \textbf{Know Your Enemy} --- Intellectual history of the American 
          right. For a Pynchon fan who wants to understand the paranoid style 
          in American politics.
\end{enumerate}

\subsection{Literary / Pynchon-Adjacent}
\begin{enumerate}[label=\textbf{\arabic*.}]
    \item \textbf{Death Is Just Around The Corner} --- 
          \emph{THE Pynchon podcast.} Deep, obsessive readings of 
          \emph{Gravity's Rainbow}, \emph{Mason \& Dixon}, etc. 
          Also covers other postmodern lit and conspiracy culture.
    \item \textbf{Entitled Opinions} (Stanford) --- A literature professor 
          discusses great books, philosophy, and ideas. Deeply intellectual.
    \item \textbf{Backlisted} --- A UK podcast about forgotten or 
          underappreciated books. Smart, warm, surprising choices.
\end{enumerate}

% =============================================================
\newpage
\section{First vs. Second Conditional --- Quick \& Dirty}
% =============================================================

\begin{tcolorbox}[colback=red!5!white, colframe=red!60!black, 
    title=\textbf{FIRST CONDITIONAL --- Real / Possible Future}]
\textbf{Structure:} \texttt{If + present simple, will + base verb}

\textbf{Use:} Something that \emph{can really happen}.

\textbf{Examples:}
\begin{itemize}
    \item If it \textbf{rains} tomorrow, I \textbf{will stay} home.
    \item If you \textbf{study} hard, you \textbf{will pass} the TOEFL.
    \item If Aldana \textbf{publishes} a new article, I \textbf{will translate} it.
\end{itemize}

\textbf{Vibe:} ``This could legit happen.''
\end{tcolorbox}

\vspace{0.5cm}

\begin{tcolorbox}[colback=blue!5!white, colframe=blue!60!black, 
    title=\textbf{SECOND CONDITIONAL --- Unreal / Hypothetical / Unlikely}]
\textbf{Structure:} \texttt{If + past simple, would + base verb}

\textbf{Use:} Something \emph{imaginary, unlikely, or contrary to reality}.

\textbf{Examples:}
\begin{itemize}
    \item If I \textbf{were} Pynchon, I \textbf{would never} do interviews.
    \item If I \textbf{spoke} perfect English, I \textbf{wouldn't need} the TOEFL.
    \item If Mexico \textbf{had} better infrastructure, everything \textbf{would change}.
\end{itemize}

\textbf{Vibe:} ``This is imaginary / not happening right now.''
\end{tcolorbox}

\vspace{0.5cm}

\begin{tcolorbox}[colback=gray!5!white, colframe=black, 
    title=\textbf{TL;DR --- The Key Difference}]
\begin{center}
\begin{tabular}{|l|l|l|}
\hline
\textbf{Conditional} & \textbf{Probability} & \textbf{Tense Trick} \\
\hline
1st (First) & Possible / Real & \texttt{If + PRESENT $\rightarrow$ WILL} \\
\hline
2nd (Second) & Unreal / Hypothetical & \texttt{If + PAST $\rightarrow$ WOULD} \\
\hline
\end{tabular}
\end{center}

\textbf{Memory hack:} \\
1st = ``If I \textbf{have} time, I \textbf{will} go to the gym.'' (maybe I will!) \\
2nd = ``If I \textbf{had} a million dollars, I \textbf{would} buy a house.'' (I don't have it!)
\end{tcolorbox}

% =============================================================
\section{Next Steps}
% =============================================================

\begin{enumerate}
    \item \textbf{Dime ``Dame otro artículo de Aldana''} y te paso uno en español 
          para que lo traduzcas al inglés. Te corrijo con feedback detallado.
    \item \textbf{Dime ``Podcast episode rec''} y te sugiero un episodio 
          específico para tu camino al gym.
    \item \textbf{Dime ``Grammar drill''} y hacemos ejercicios de condicionales 
          u otro tema.
    \item Todo lo seguimos haciendo en \LaTeX.
\end{enumerate}

\end{document}